% !TeX program = lualatex
\documentclass[12pt,a4paper]{article}

% ========= حزمة =========
\usepackage[english, bidi=default, layout=counters]{babel}
\usepackage{amsmath,amssymb,amsthm,mathtools}
\usepackage{geometry}
\geometry{left=1.0cm,right=1.0cm,top=1.25cm,bottom=3.1cm}
\usepackage{booktabs}
\usepackage{array}
\usepackage{hyperref}
\usepackage{url}
\usepackage{graphicx}
\usepackage{caption}
\usepackage{subcaption}
\usepackage{float}
\usepackage{siunitx}
\usepackage{bm}
\usepackage{microtype}
\usepackage{fancyhdr}
\usepackage{fontspec}
\usepackage{parskip}
\usepackage{xcolor}
\usepackage{tikz}
\usepackage{pgfplots}
\pgfplotsset{compat=1.18}
\usetikzlibrary{arrows.meta, positioning, calc, shapes.geometric}
\usepackage{nicefrac}
\usepackage{enumitem}
\usepackage{makecell}
\usepackage{ragged2e}
\usepackage{tabularx}
\usepackage{slashed}

% ========= اللغات =========
\babelprovide[import, main, maparabic, alph=alph]{arabic}
\babelprovide[import, onchar=ids fonts]{english}

% ========= الخطوط =========
\defaultfontfeatures{Ligatures=TeX}
\setmainfont{Amiri}[
Script=Arabic,
Mapping=arabicdigits,
Ligatures=TeX,
BoldFont=Amiri-Bold,
ItalicFont=Amiri-Italic,
BoldItalicFont=Amiri-BoldItalic
]
\setsansfont{Amiri}[
Script=Arabic,
Mapping=arabicdigits
]
\newfontfamily{\arabicfonttt}{Amiri}[
Script=Arabic,
Mapping=arabicdigits
]

% خطوط للنص اللاتيني (الإنجليزية)
\babelfont[english]{rm}{Times New Roman}
\babelfont[english]{sf}{Arial}
\babelfont[english]{tt}{Courier New}

% ========= ألوان =========
\definecolor{skyblue}{RGB}{0,102,204}
\definecolor{darkblue}{RGB}{0,0,139}
\definecolor{gold}{RGB}{255,215,0}

% ========= أوامر YUCT =========
\newcommand{\Keff}{K_{\mathrm{eff}}}
\newcommand{\Keffzero}{K_{\mathrm{eff},0}}
\newcommand{\kappac}{\kappa_c}
\newcommand{\eps}{\varepsilon}
\newcommand{\df}{d_{\!f}}
\newcommand{\constmin}{C_{\min}}
\newcommand{\dint}{\ensuremath{D\!+\!I\!\cdot\!R}}
\newcommand{\Mpl}{M_{\mathrm{Pl}}}
\newcommand{\mpl}{m_{\mathrm{Pl}}}
\newcommand{\Lpl}{l_{\mathrm{Pl}}}
\newcommand{\RH}{R_{\mathrm{H}}}
\newcommand{\tH}{t_{\mathrm{H}}}
\newcommand{\ninf}{N_{\mathrm{e}}}
\newcommand{\ns}{n_{\mathrm{s}}}
\newcommand{\nt}{n_{\mathrm{T}}}
\newcommand{\rs}{r}
\newcommand{\alD}{\alpha_{\mathrm{D}}}
\newcommand{\beI}{\beta_{\mathrm{I}}}
\newcommand{\gaR}{\gamma_{\mathrm{R}}}
\newcommand{\Kcrit}{K_{\mathrm{crit}}}
\newcommand{\Rstruct}{R_{\mathrm{struct}}}
\newcommand{\Ntrials}{N_{\mathrm{trials}}}
\newcommand{\Nlife}{\langle N_{\mathrm{life}}\rangle}

% ========= عوامل =========
\DeclareMathOperator{\Tr}{Tr}
\DeclareMathOperator{\Var}{Var}
\DeclareMathOperator{\Cov}{Cov}
\DeclareMathOperator{\Div}{Div}
\DeclareMathOperator{\Grad}{Grad}

% ========= نظريات =========
\newtheorem{theorem}{نظرية}[section]
\newtheorem{lemma}[theorem]{ليما}
\newtheorem{definition}[theorem]{تعريف}
\newtheorem{corollary}[theorem]{نتيجة}
\newtheorem{proposition}[theorem]{اقتراح}
\theoremstyle{remark}
\newtheorem{remark}[theorem]{ملاحظة}
\newtheorem{example}[theorem]{مثال}

% ========= رؤوس وتذييلات =========
\fancypagestyle{firstpage}{%
	\fancyhf{}%
	\renewcommand{\headrulewidth}{0pt}%
	\renewcommand{\footrulewidth}{0pt}%
	\fancyfoot[C]{%
		\begin{minipage}[t]{\textwidth}
			\centering
			\textcolor{gold}{\rule{\textwidth}{1.6pt}}\\[5pt]
			{\small \textsuperscript{1}أبحاث ياكوشيف، \url{https://yuct.org/}} \hfill
			{\small بروتوكول YPSDC: \url{https://ypsdc.com/}}\\[-2pt]
			\textcolor{gold}{\rule{\textwidth}{1.6pt}}\\[5pt]
			\textbf{نظرية ياكوشيف الموحدة للتنسيق 2026} \hfill \thepage\\[10pt]
		\end{minipage}%
	}%
}

\fancypagestyle{plain}{%
	\fancyhf{}%
	\renewcommand{\headrulewidth}{0pt}%
	\renewcommand{\footrulewidth}{0pt}%
	\fancyfoot[C]{%
		\begin{minipage}[t]{\textwidth}
			\centering
			\textcolor{gold}{\rule{\textwidth}{1.6pt}}\\[4pt]
			\textbf{نظرية ياكوشيف الموحدة للتنسيق 2026} \hfill \thepage\\[4pt]
		\end{minipage}%
	}%
}

\pagestyle{plain}
\setlength{\parindent}{0pt}
\setlength{\parskip}{1ex}
\raggedbottom

\hypersetup{
	colorlinks=true,
	linkcolor=blue,
	citecolor=blue,
	urlcolor=blue,
}

% ========= رأس المقال (شعار YUCT) =========
\newcommand{\makeheader}{%
	\noindent
	\begin{minipage}{0.17\textwidth}
		\raggedright
		{\sffamily\bfseries\color{skyblue}\fontsize{46}{46}\selectfont YUCT}%
	\end{minipage}%
	\hspace{30pt}
	\begin{minipage}{0.32\textwidth}
		\raggedright
		{\sffamily\fontsize{11}{13}\selectfont نظرية ياكوشيف\\ الموحدة\\ للتنسيق}%
	\end{minipage}%
	\hfill
	\begin{minipage}{0.38\textwidth}
		\raggedleft
		{\sffamily\bfseries ملحق W}\\
		{\sffamily مارس 2026}\\
		{\sffamily\footnotesize \scriptsize DOI: \url{https://doi.org/10.5281/zenodo.18444598}}%
	\end{minipage}
	\par\vspace{5pt}
	\noindent\textcolor{gold}{\rule{\textwidth}{1.6pt}}
	\par\vspace{0pt}
}

\begin{document}
	\thispagestyle{firstpage}
	\makeheader
	
	\vspace{0cm}
	{\LARGE\bfseries التصوير المقطعي الجذبي السلبي للباطن الكوكبي باستخدام الأمواج المدية \\
		\Large مقاربة قائمة على التنسيق ضمن نظرية ياكوشيف الموحدة للتنسيق (YUCT) \par}
	\vspace{0.2cm}
	
	{\large أليكسي ف. ياكوشيف\footnote{أبحاث ياكوشيف، \url{https://yuct.org/}}} \\
	\vspace{0.0cm}
	
	{\color{darkblue}\textbf{\LARGE المستخلص}} \par
	{\large
		\noindent
		يقدم هذا الملحق طريقة جديدة للاستكشاف الجيوفيزيائي السلبي قائمة على تحليل الاستجابات القشرية الرنينية للتأثير الجذبي المدّي من القمر والشمس. ضمن نظرية ياكوشيف الموحدة للتنسيق (YUCT)، نظهر أن طيف التشوهات المدية يحتوي على معلومات حول التوزيع تحت السطحي لكفاءة التنسيق \(\Keff\)، المرتبطة مباشرة بنوع الصخر، والمحتوى المائع (نفط، غاز، ماء)، والتمعدنات الخام. نطور نموذجاً رياضياً يربط قياسات الجاذبية والزلازل السطحية بتوزيع \(\Keff\) العمقي، مستفيدين من قانون قياس الخطأ الشامل \(\eps = \kappac \alpha \Keff^{-\beta}\) (\(\beta=0.67\), \(\kappac=0.33\)) ومعادلات الاقتران بين حقول الجاذبية والتشوهات المرنة من قطاعي YUCT 0 (الجاذبية) و34 (علم الكواكب). نظهر أن الأمواج المدية تقدم طريقة سبر غير مكلفة، ومفيدة بيئياً، وقابلة للتطبيق عالمياً، بأعماق اختراق لا تبلغها المصادر النشطة. نستعرض أبحاثاً مستقلة تستبق عناصر من هذه المقاربة، بما في ذلك أرصاد مباشرة للرنين المدّي ومقترحات لاستكشاف الهيدروكربونات باستخدام المد الجذبي. تُقترح منهجية لتحديد الهيدروكربونات عبر ترددات رنينية مميزة وعوامل جودة، بناءً على نظام تصنيف معادن وموائع قائم على التنسيق مماثل للجدول الدوري للعناصر. تُقيم الجدوى الاقتصادية: تخفيض خطر الحفر الجاف بنسبة 20–30\% سيسترجع تكاليف النشر عدة مرات. نوسع المنهجية أيضاً إلى مراقبة الصحة البنيوية للمنشآت الهندسية (ناطحات سحاب، جسور، سدود)، والتنبؤ بالزلازل، والتنبؤ بأمواج تسونامي، مبينين أن نفس المبادئ الفيزيائية تمكن من فئة ثورية جديدة من تقنيات المراقبة السلبي ذات القدرة على إنقاذ ملايين الأرواح وتريليونات الدولارات. تفتح النتائج طريقاً إلى فئة جديدة جوهرياً من تقنيات الاستكشاف والمراقبة باستخدام حقول الجاذبية الطبيعية.
	}
	
	\vspace{0.1cm}
	\noindent
	{\color{darkblue}\textbf{الكلمات المفتاحية:}} YUCT، الأمواج المدية، كفاءة التنسيق، تصوير مقطعي عميق، استكشاف الهيدروكربونات، طريقة رنينية، جيوفيزياء، باطن كوكبي، مراقبة الصحة البنيوية، تنبؤ بالزلازل، تنبؤ بتسونامي.
	
	\vspace{0.5cm}
	\noindent
	{\small نُسخة مختصرة من هذا العمل نُشرت سابقاً ومتاحة على \\ \url{https://doi.org/10.5281/zenodo.18362308}}
	\vspace{0.5cm}
	
	\noindent
	\textcopyright\ 2026 أبحاث ياكوشيف. جميع الحقوق محفوظة.
	
	\tableofcontents
	\newpage
	
	\section{مقدمة: تحدي الجيوفيزياء العميقة}
	\label{sec:intro}
	
	\subsection{حدود الاستكشاف النشط}
	
	يبقى استكشاف رواسب الهيدروكربونات والمعادن أحد أكثر المساعي كثافة رأس مالية في علوم الأرض التطبيقية. تتطلب طرق الانعكاس الزلزالي التقليدية مصادر اصطناعية قوية (متفجرات، شاحنات Vibroseis، مدافع هوائية)، مما يحد من أعماق الاختراق إلى 5–10 كم ويتكبد تكاليف بيئية كبيرة. تقدم مسوحات الجاذبية والمغناطيسية فقط معلومات متكاملة عمقياً ولا يمكنها تحديد أنواع الموائع بشكل لا لبس فيه.
	
	\subsection{المد والجزر كإشارة سبر طبيعية}
	
	تتشوه الأرض باستمرار بفعل قوى الجاذبية من القمر والشمس، منتجة مداً أرضياً صلباً بسعات تصل إلى 30–50 سم عند خط الاستواء. هذه التشوهات دورية تماماً، بطيف توافقي غني يمتد عبر مركبات نصف يومية (12.42 س، 12.00 س)، ويومية (24–26 س)، وطويلة الدور (14 ي، 6 أشهر، 18.6 سنة). تقليدياً، تُعامل الإشارات المدية كضجيج يجب إزالته من القياسات عالية الدقة.
	
	يقترح هذا الملحق تحولاً في النموذج: \textbf{استخدام الأمواج المدية كإشارة سبر طبيعية سلبية} لسبر باطن الأرض. ضمن نظرية ياكوشيف الموحدة للتنسيق (YUCT)، نظهر أن الاستجابة المدية تشفر معلومات حول التوزيع تحت السطحي لكفاءة التنسيق \(\Keff\)، التي تعمل كمؤشر شامل لخواص الصخور، والمحتوى المائع، والتمعدن الخام.
	
	\subsection{بنية هذا الملحق}
	
	يستعرض القسم \ref{sec:prior} أبحاثاً مستقلة تستبق عناصر من هذه المقاربة، مؤسساً أسسها التجريبية. يستذكر القسم \ref{sec:yuct-foundations} مبادئ YUCT ذات الصلة. يطور القسم \ref{sec:model} النموذج الرياضي الرابط بين التشوهات المدية و\(\Keff\)، محدداً صراحة قطاعات YUCT المساهمة وحدود الاقتران في اللاغرانجيان. يقدم القسم \ref{sec:identification} نظام تصنيف قائم على التنسيق للمواد الجيولوجية. يقيم القسم \ref{sec:economics} الجدوى الاقتصادية. يوجز القسم \ref{sec:roadmap} خريطة طريق تجريبية. يناقش القسم \ref{sec:conclusion} الآثار على استكشاف الكواكب والمهمات المستقبلية.
	
	\subsection{من استكشاف الكواكب إلى مراقبة الصحة البنيوية}
	\label{sec:beyond-geophysics}
	
	للمنهجية المطورة في هذا الملحق لسبر الباطن الكوكبي امتداد طبيعي وقوي: التقييم غير المدمر للمنشآت التي صنعها الإنسان. أي جسم ضخم — ناطحة سحاب، جسر، سد، وعاء احتواء مفاعل نووي — يتشوه باستمرار بنفس قوى المد الجذبية التي تشكل الأجسام الكوكبية. هذه التشوهات، رغم صغرها، تحمل في طيفها بصمة الحالة الداخلية للجسم: توزيع صلابته، وجود تشققات، تدهور المواد، وكفاءة التنسيق الكلية \(K_{\mathrm{eff}}\) للمنشأ.
	
	يحول هذا التبصر التصوير المقطعي المدّي من أداة جيوفيزيائية بحتة إلى طريقة شاملة لـ\textbf{مراقبة الصحة البنيوية السلبية}. على عكس الطرق النشطة التي تتطلب إثارة اصطناعية (هزازات، صدمات، فوق صوتيات)، المراقبة المدية هي:
	\begin{itemize}
		\item \textbf{سلبية:} تستخدم إشارات جذبية طبيعية، مجانية، ومستمرة.
		\item \textbf{شاملة:} تقدم صورة متكاملة للمنشأ بأكمله، وليس فقط نقاطاً محلية.
		\item \textbf{عميقة:} تخترق الحجم بأكمله، كاشفة عن ضرر داخلي خفي.
		\item \textbf{آمنة:} لا إشعاع، لا مدخلات عالية الطاقة، لا حاجة للتلامس.
		\item \textbf{مستمرة:} تمكن من تقييم في الزمن الحقيقي أو دوري دون انقطاع الخدمة.
	\end{itemize}
	
	العلاقة الأساسية \(Q \propto K_{\mathrm{eff}}^{\beta}\) (المعادلة \ref{eq:Q-keff}) تربط عامل الجودة القابل للرصد لرنين بنيوي مباشرة بكفاءة تنسيق المادة. مع تقدم المنشأ في العمر، وتراكم التشققات المجهرية، أو تعرضه للضرر، تنخفض \(K_{\mathrm{eff}}\) لديه، وينخفض عامل الجودة لرنيناته المدية تبعاً لذلك. مراقبة هذه التغيرات عبر الزمن تقدم مقياساً كمياً غير باضع للصحة.
	
	\section{استباقات مستقلة وأسس تجريبية}
	\label{sec:prior}
	
	\subsection{أرصاد مباشرة للرنين المدّي}
	
	بشكل لافت، تلقت الفكرة الأساسية لهذا الملحق — أن الأمواج المدية يمكنها التفاعل رنينياً مع البنى الجيولوجية — تأكيداً رصدياً مباشراً بالفعل.
	
	\begin{quotation}
		\noindent
		\textit{"سُجلت رنينات المد الجذبي طويل الدور (14–15 يوماً) وأمواج التشوه في القشرة الأرضية التي حدثت بسبب إعادة تموضع مركز كتلة نظام الأرض-قمر في مناطق نشطة زلزالياً من ألطاي-سايان، وكامتشاتكا، وسخالين أثناء عملية رصد الزلازل... تثبت المقالة أن رنينات المد الجذبي يمكن استخدامها للتنبؤ بالزلازل وكذلك للإنذار المبكر بالمخاطر الجيوديناميكية للمنشآت الهيدروليكية \textbf{وللاستكشاف المباشر لخزانات النفط والغاز}."} \cite{Sibgatulin2014}
	\end{quotation}
	
	يقدم عمل Sibgatulin et al. \cite{Sibgatulin2014} عدة أسس تجريبية حاسمة للنظرية الحالية:
	
	\begin{itemize}
		\item \textbf{رصد الرنينات:} تنتج مركبات المد طويل الدور (14–15 يوماً) استجابات رنينية قابلة للقياس في القشرة.
		\item \textbf{آلية التحفيز:} تعمل هذه الرنينات كمحفزات لزلازل قوية (M\(\ge\)5.0) في 80\% من الحالات، مؤكدة أن الطاقة المدية تقترن بشكل فعال بالبنى القشرية.
		\item \textbf{أمواج التشوه:} تولد الرنينات أمواج تشوه بطيئة (1–5 كم/سا)، مفسرة كسوليتونات، تنتشر عبر القشرة.
		\item \textbf{إمكانات استكشافية:} يقترح المؤلفون صراحة استخدام هذه الرنينات للاستكشاف المباشر للهيدروكربونات — استباق ثاقب للمنهجية المطورة هنا.
	\end{itemize}
	
	\subsection{التصوير المقطعي المدّي في علم الكواكب}
	
	دخل مصطلح "التصوير المقطعي المدّي" مؤخراً أدبيات علم الكواكب. يوضح Rovira-Navarro et al. \cite{RoviraNavarro2025} أن:
	
	\begin{itemize}
		\item ستقيس تجربة 3GM Radio Science Experiment لمهمة JUICE تغيرات الجاذبية المحدثة مدياً على جانيميد بدقة كافية لرسم خرائط التغيرات الجانبية في البنية الداخلية.
		\item تنتج تغيرات سمك القشرة بنسبة 10–100\% إشارات مدية قابلة للكشف، مع أرصاد توافقية كروية من الدرجة 2 والدرجة 3 تقيد فروق السمك بين خط الاستواء والقطب وثنائيات نصف الكرة.
		\item يذكر المؤلفون صراحة أن "التصوير المقطعي المدّي" — تقنية استُخدمت سابقاً على الأرض — يمكن توسيعها لتشمل الأقمار الجليدية.
	\end{itemize}
	
	يدعم هذا العمل المبدأ الأساسي بأن الاستجابات المدية حساسة للبنية الداخلية ثلاثية الأبعاد، داعماً المقاربة المطورة هنا للأرض.
	
	\subsection{قياسات مدية تحت سطحية}
	
	أجرى Rogister et al. \cite{Rogister2024} تجربة فريدة بمقارنة إشارات الجاذبية المدية عند السطح وعند عمق 520 م في مختبر Rustrel تحت الأرض، فرنسا. بينما كان تركيزهم الأساسي على معايرة مقاييس الجاذبية، يوضح عملهم أن:
	
	\begin{itemize}
		\item يمكن قياس تغيرات الجاذبية المدية آنياً عند أعماق مختلفة بمقاييس جاذبية فائقة الموصلية.
		\item الفرق النظري بين الإشارات المدية السطحية وتحت السطحية صغير (\(8.5 \times 10^{-5}\) للأمواج نصف اليومية) لكنه قابل للكشف من حيث المبدأ.
		\item تحد شكوك المعايرة الحالية من القدرة على حل هذا الفرق، لكن تحسين الآليات يمكن أن يمكن من قياسات مدية معتمدة على العمق.
	\end{itemize}
	
	\subsection{الاستجابة المدية لعطارد وكواكب أخرى}
	
	يدرس Briaud et al. \cite{Briaud2025} الاستجابة المدية لعطارد باستخدام بيانات MESSENGER ويستعدون لأرصاد BepiColombo. يشددون على أن:
	
	\begin{itemize}
		\item أعداد لوف المدية حساسة للتباينات الناشئة عن تغيرات مكانية في درجة الحرارة، والتركيب، والخواص الفيزيائية.
		\item بشكل مماثل للأقمار الجليدية، تغير الفروق الحرارية أو المعدنية الموضعية الاستجابات الميكانيكية، مما يؤدي إلى تشوه مدّي غير منتظم.
		\item قياس الميسان طويل الدور، وتأخر الطور المدّي، والانحرافات عن حالة كاسيني يمكن أن تقيد البنية الداخلية.
	\end{itemize}
	
	تظهر هذه التطبيقات الكوكبية كونية التصوير المقطعي المدّي كأداة لسبر البنية الداخلية عبر أجسام متنوعة.
	
	\subsection{توليف}
	
	يؤسس البحث المستقل المستعرض أعلاه أن:
	
	\begin{enumerate}
		\item \textbf{الرنينات المدية قابلة للرصد} وترتبط بالبنية الجيولوجية \cite{Sibgatulin2014}.
		\item \textbf{التصوير المقطعي المدّي تقنية معترف بها} في علم الكواكب لرسم خرائط البنية الداخلية ثلاثية الأبعاد \cite{RoviraNavarro2025}.
		\item \textbf{القياسات المدية تحت السطحية} ممكنة تقنياً بالآليات الحالية \cite{Rogister2024}.
		\item \textbf{الاستجابات المدية الكوكبية} حساسة للتباينات الداخلية \cite{Briaud2025}.
	\end{enumerate}
	
	ما كان مفقوداً هو إطار نظري موحد يربط هذه الأرصاد بمعامل كمي (\(\Keff\)) يمكن عكسه لخواص تحت السطح. هذا بالضبط ما تقدمه YUCT.
	
	\section{أسس YUCT للتصوير المقطعي المدّي}
	\label{sec:yuct-foundations}
	
	\subsection{كفاءة التنسيق \(\Keff\) وقانون الخطأ الشامل}
	
	ضمن YUCT، يتميز أي نظام فيزيائي بكفاءة تنسيقه \(\Keff\)، التي تحدد درجة الترابط والتزامن بين مكوناته. يربط قانون قياس الخطأ الشامل (الملحق L) أي خطأ نسبي أو تذبذب \(\eps\) بـ \(\Keff\):
	
	\begin{equation}
		\eps = \kappac \, \alpha \, \Keff^{-\beta}, \qquad \beta = 0.67 \pm 0.03,\ \kappac = 0.33,
		\label{eq:universal-scaling}
	\end{equation}
	
	حيث \(\alpha\) ثابت يعتمد على المنظومة من مرتبة \(0.01\)–\(1\). تم التحقق من هذا القانون عبر 40 مرتبة أسية، من أخطاء تضاعف DNA إلى تذبذبات الخلفية الكونية الميكروية.
	
	في سياق التصوير المقطعي المدّي، يمثل \(\eps\) التوهين النسبي أو تأخر الطور للتشوهات المدية، بينما تميز \(\Keff\) الوسط الجيولوجي (نوع الصخر، المحتوى المائع، المسامية، التشقق).
	
	\subsection{بنية القطاعات في YUCT ذات الصلة بالتصوير المقطعي المدّي}
	
	يقسم لاغرانجيان YUCT (الملحق A) الواقع إلى 120 قطاعاً متفاعلاً. للتصوير المقطعي المدّي، القطاعات التالية ذات صلة مباشرة:
	
	\begin{itemize}
		\item \textbf{القطاع 0 (الجاذبية):} يصف المترية \(g_{\mu\nu}\) والحقل الجذبي. الكمون المدّي \(W(\mathbf{r},t)\) المولد من القمر والشمس يدخل عبر هذا القطاع.
		\item \textbf{القطاع 34 (علم الكواكب):} يصف البنية الداخلية للكواكب، بما في ذلك الكثافة \(\rho(\mathbf{r})\)، وثوابت لامي \(\lambda(\mathbf{r}), \mu(\mathbf{r})\)، والتشوهات المدية.
		\item \textbf{القطاع 1 (الكهرومغناطيسية):} ذو صلة بالارتباطات المغناطيسوتلورية والسلائف الكهرومغناطيسية للرنينات المدية (كما رصدها Sibgatulin et al. \cite{Sibgatulin2014} عبر قياسات الحقل الكهرومغناطيسي النبضي الطبيعي).
		\item \textbf{القطاع 62 (البيئة) و86 (الديموغرافيا):} ذو صلة لتقييم الآثار البيئية والاجتماعية-الاقتصادية لأنشطة الاستكشاف (القسم \ref{sec:economics}).
	\end{itemize}
	
	يتم الاقتران بين القطاعات بوساطة معاملات \(\kappa_{sr}\). للتصوير المقطعي المدّي، الحد الحرج هو:
	
	\begin{equation}
		\mathcal{L}_{\text{اقتران}} = \kappa_{0,34} \, \mathrm{Tr}\left(\Psi_{0,34} \cdot O_0 \cdot O_{34}^\dagger\right),
		\label{eq:coupling}
	\end{equation}
	
	حيث \(\Psi_{0,34}\) هو حقل التنسيق الرابط بين الجاذبية والبنية الكوكبية، و\(O_0\) يمثل مؤثر الكمون المدّي، و\(O_{34}\) يمثل مؤثر موتر التشوه.
	
	\subsection{اعتماد درجة الحرارة والبنية العميقة}
	
	يوضح الملحق P أن \(\Keff\) تعتمد على درجة الحرارة، مع \(\Keff(T) \propto T^{-\gamma}\) و\(\beta\gamma = 0.20\) المحددين من البيانات الجيوفيزيائية (نظام الأرض-قمر، انزياحات الأقزام البيضاء نحو الأحمر). هذه الحساسية الحرارية حاسمة للتصوير المقطعي العميق، حيث تقدم الممالات الحرارية الأرضية قيوداً إضافية على مقاطع \(K_{\mathrm{eff}}\). علاوة على ذلك، فهي تعني أن \textbf{التضمين الحراري للرنينات المدية} يمكن استغلاله لفصل المساهمات من أعماق مختلفة، لأن تغيرات درجة حرارة السطح تنتشر انتشارياً وتؤثر على الطبقات الأقل عمقاً بسرعة أكبر (انظر القسم~\ref{sec:thermal-modulation}).
	
	\section{النموذج الرياضي للاستجابة المدية}
	\label{sec:model}
	
	\subsection{التشوه المرن تحت الكمون المدّي}
	
	اعتبر نموذج أرض مرن بكثافة \(\rho(\mathbf{r})\)، ومعاملات لامي \(\lambda(\mathbf{r}), \mu(\mathbf{r})\)، خاضع لكمون مدّي \(W(\mathbf{r},t)\) من أجسام خارجية. معادلة الحركة للإزاحات الصغيرة \(\mathbf{u}(\mathbf{r},t)\) هي:
	
	\begin{equation}
		\rho \frac{\partial^2 \mathbf{u}}{\partial t^2} = \nabla \cdot \boldsymbol{\sigma} + \rho \nabla W,
		\label{eq:motion}
	\end{equation}
	
	حيث \(\boldsymbol{\sigma}\) هو موتر الإجهاد. لوسط موحد الخواص مرن:
	
	\begin{equation}
		\sigma_{ij} = \lambda \delta_{ij} \nabla \cdot \mathbf{u} + \mu \left(\frac{\partial u_i}{\partial x_j} + \frac{\partial u_j}{\partial x_i}\right).
		\label{eq:stress}
	\end{equation}
	
	يُفرد الكمون المدّي بتوافقيات كروية:
	
	\begin{equation}
		W(r,\theta,\phi,t) = \sum_{n=2}^{\infty} \sum_{m=-n}^{n} W_{nm}(r) Y_{nm}(\theta,\phi) e^{i\omega_{nm}t}.
		\label{eq:tidal-potential}
	\end{equation}
	
	للأرض الصلبة، تهيمن توافقيات الدرجة 2 (الأمواج نصف اليومية واليومية). يُعبر عن الحل بنشر توافقي كروي بدوال قطرية تحقق نظام معادلات لوف.
	
	\subsection{إدماج \(\Keff\) في المعاملات المرنة}
	
	في YUCT، يُعبر عن الخواص المرنة عبر \(\Keff\). باتباع منهجية الملحق C (كيمياء التنسيق)، نقدم \textbf{نظام تصنيف معادن قائم على التنسيق}. لصخر أو مائع معطى، ترتبط \(\Keff\) بمعامل الحجم \(K\)، ومعامل القص \(\mu\)، والكثافة \(\rho\) بعلاقات معايرة تجريبياً:
	
	\begin{align}
		\lambda(\mathbf{r}) &= \lambda_0 \cdot f_\lambda(\Keff(\mathbf{r})), \\
		\mu(\mathbf{r}) &= \mu_0 \cdot f_\mu(\Keff(\mathbf{r})), \\
		\rho(\mathbf{r}) &= \rho_0 \left(\frac{\Keff(\mathbf{r})}{\Keffzero}\right)^{\gamma_\rho}, \quad \gamma_\rho \approx 0.2\ \text{(مدى 0.1–0.3)},
		\label{eq:keff-properties}
	\end{align}
	
	الدوال \(f_\lambda, f_\mu\) تزايدية رتيبة، معكسة أن التنسيق الأعلى (بنية ذرية/جزيئية أكثر ترابطاً) يعطي صلابة أكبر. تقترح التقديرات الأولية \(f_\lambda(\Keff) \propto \Keff^{\xi}\) مع \(\xi \approx 0.3\)–\(0.5\)، استناداً إلى قياس المعاملات المرنة مع الكثافة وعدد التنسيق.
	
	\subsection{الاستجابة الرنينية وعامل الجودة}
	
	عندما يقترب التردد \(\omega\) لتوافقي مدّي من التردد الطبيعي \(\omega_0\) لبنية تحت سطحية (مثلاً خزان مملوء بالموائع)، يحدث تضخيم رنيني. يرتبط عامل الجودة \(Q\) للرنين بـ \(\Keff\) عبر قانون الخطأ الشامل. تبديد الطاقة لكل دورة يتناسب مع الخطأ النسبي \(\eps\)، وبما أن \(1/Q = \Delta E / (2\pi E) \approx \eps\)، لدينا:
	
	\begin{equation}
		\frac{1}{Q} \propto \eps = \kappac \alpha \Keff^{-\beta} \quad \Longrightarrow \quad Q = Q_0 \, \Keff^{\beta},
		\label{eq:Q-keff}
	\end{equation}
	
	حيث \(Q_0\) قيمة مرجعية. وهكذا:
	\begin{itemize}
		\item الأوساط عالية \(\Keff\) (صخور متماسكة، موائع كثيفة) تنتج رنينات حادة وعالية \(Q\).
		\item الأوساط منخفضة \(\Keff\) (مناطق متشققة، غاز) تنتج سمات عريضة ومنخفضة \(Q\).
	\end{itemize}
	
	\subsection{النموذج الأمامي والمسألة العكسية}
	
	تنتج القياسات السطحية لتغيرات الجاذبية والإزاحات سلاسل زمنية \(d(t)\). بعد استخراج التوافقيات المدية (عبر التفاف مع إفيميريدات معروفة) وطرح استجابة نموذج مرجعي متناظر كروياً، يحتوي المتبقي \(\delta d(t)\) على معلومات حول التباينات الجانبية.
	
	يعطي تحويل فورييه الطيف \(\delta D(\omega)\)، المحتوي على قمم عند ترددات \(\omega_i\) المقابلة لرنينات البنى تحت السطحية. من كل قمة نستخرج:
	
	\begin{itemize}
		\item \(\omega_i\) — التردد الرنيني، حساس للعمق، والحجم، والخواص المرنة.
		\item \(Q_i = \omega_i / \Delta \omega_i\) — عامل الجودة، مرتبط بـ \(\Keff\) عبر المعادلة (\ref{eq:Q-keff}).
		\item \(A_i\) — السعة، متناسبة مع تباين \(\Keff\) بين البنية والخلفية.
	\end{itemize}
	
	تُحل المسألة العكسية — استعادة \(\Keff(\mathbf{r})\) ثلاثي الأبعاد — عبر عكس تصويري مقطعي، بتصغير عدم التطابق بين الأطياف المرصودة والاصطناعية. يدمج التنظيم معلومات مسبقة من نظام التصنيف المعدني القائم على التنسيق.
	
	\subsection{التضمين الحراري للرنينات المدية}
	\label{sec:thermal-modulation}
	
	تؤثر تغيرات درجة الحرارة على \(\Keff\) عبر العلاقة \(\Keff(T) \propto T^{-\gamma}\) (الملحق P). تنتشر تغيرات درجة الحرارة اليومية والموسمية في باطن الأرض كأمواج حرارية، مضمنة \(\Keff\) المحلية وبالتالي الترددات الرنينية وعوامل الجودة للبنى الضحلة. يقدم هذا بعداً إضافياً لتمييز العمق: تستجيب البنى الأعمق بتأخر طور وسعة مخفضة لتأثير درجة حرارة السطح.
	
	إذا كان حقل درجة الحرارة \(T(\mathbf{r},t)\) معروفاً من بيانات الأرصاد الجوية أو نموذجاً، يمكن حساب الاضطراب المعتمد على الزمن \(\delta \Keff(\mathbf{r},t) \approx -\gamma \Keff(\mathbf{r}) \delta T(\mathbf{r},t)/T\). يضمن هذا الاستجابة المدية، منتجاً نطاقات جانبية أو تضمين سعة عند ترددات حرارية. تحليل هذه التضمينات يسمح بفصل المساهمات الضحلة والعميقة، محسناً الاستبانة العمودية.
	
	\subsection{الامتداد إلى المنشآت الهندسية: ناطحات السحاب، والجسور، والسدود}
	\label{sec:structures}
	
	ينطبق النموذج الرياضي المطور في الأقسام \ref{sec:model} مباشرة على المنشآت البشرية الكبيرة، مع تعديلات مناسبة على الشروط الحدية والهندسة. لمنشأ كتلته \(M\)، وارتفاعه \(H\)، وصلابته المميزة \(k\)، يتدرج تردد النمط الانعطافي الأساسي كـ:
	
	\begin{equation}
		f_0 \approx \frac{1}{2\pi} \sqrt{\frac{k}{M_{\mathrm{eff}}}} \sim \frac{1}{H} \sqrt{\frac{E}{\rho}},
		\label{eq:structural-frequency}
	\end{equation}
	
	حيث \(E\) هو معامل يونغ، \(\rho\) هي الكثافة، و\(M_{\mathrm{eff}}\) هي الكتلة النمطية الفعالة. لناطحة سحاب نموذجية (\(H \sim 300\) م، \(E/\rho \sim 10^7\) م\(^2\)/ثا\(^2\))، \(f_0 \sim 0.1\) هرتز — وهو ضمن مدى التوافقيات المدية بعد المزج غير الخطي أو الإثارة الوسيطية.
	
	عامل الجودة \(Q\) لهذه الأنماط محكوم بنفس العلاقة الشاملة:
	
	\begin{equation}
		Q = Q_0 \left( \frac{K_{\mathrm{eff}}}{K_{\mathrm{eff},0}} \right)^{\beta},
		\label{eq:structural-Q}
	\end{equation}
	
	حيث تمثل \(K_{\mathrm{eff}}\) الآن \textbf{كفاءة التنسيق البنيوية} — مقياس لابعدي لمدى نجاح جميع مكونات المنشأ في العمل معاً. المباني حديثة البناء لها \(K_{\mathrm{eff}}\) عالية؛ المنشآت المتضررة أو المتدهورة تظهر قيماً أدنى.
	
	\subsubsection{ما تكشفه المراقبة المدية}
	
	بتركيب مقاييس تسارع أو جاذبية عالية الحساسية على منشأ (مثلاً على السطح، وعند منتصف الارتفاع، وعند الأساس) وتسجيل لعدة أسابيع، نحصل على:
	
	\begin{itemize}
		\item \textbf{ترددات الأنماط:} تشير الإزاحات إلى تغيرات في الصلابة (مثلاً من التشقق أو تدهور المواد).
		\item \textbf{عوامل الجودة:} تشير الانخفاضات إلى تراكم ضرر، أو تشققات مجهرية، أو وصلات مرتخية.
		\item \textbf{تباين الخواص:} تكشف الاختلافات في الاستجابة على طول محاور مختلفة أنماط ضرر اتجاهية.
		\item \textbf{الأنماط الأعلى:} تقدم معلومات محللة عمقياً حول موقع الضرر.
	\end{itemize}
	
	\subsubsection{مفهوم "البصمة الجذبية"}
	
	لكل منشأ كبير "بصمة جذبية" فريدة — طيفه من الترددات الرنينية وعوامل الجودة المثارة بالتأثير المدّي. يمكن لهذه البصمة أن تكون:
	\begin{enumerate}
		\item \textbf{مسجلة عند التشغيل} لتأسيس خط أساس.
		\item \textbf{مراقبة دورياً} لتتبع التدهور.
		\item \textbf{معاد قياسها بعد الأحداث القصوى} (زلازل، أعاصير، حرائق) لتقييم الضرر دون تفتيش بصري.
	\end{enumerate}
	
	الحساسية لافتة: تغير 1\% في الصلابة (مثلاً من تشقق كبير) يزيح الترددات بمقدار \(\sim 0.5\%\)، وهو قابل للكشف بسهولة بالآليات الحديثة.
	
	\subsection{التنبؤ بالزلازل: مراقبة الفوالق الحرجة بالتصوير المقطعي المدّي}
	\label{sec:earthquakes}
	
	نفس قوى المد التي تمكن من تصوير الباطن الكوكبي والمنشآت الهندسية تجهد أيضاً قشرة الأرض باستمرار. لعقود، لاحظ علماء الزلازل ارتباطات إحصائية بين طور المد وحدوث الزلازل، خاصة لأحداث الفوالق العكسية والعادية \cite{Tanaka2012, Tolstoy2013, Sibgatulin2014, RAN2025}. هذه الارتباطات، رغم كونها مهمة (تصل إلى 10\% زيادة في الاحتمال عند مسافات قمرية معينة)، بقيت وصفية — ظاهرة بدون نظرية تنبؤية.
	
	تقدم YUCT الإطار المفقود. فالق جيولوجي يقترب من الانهيار هو نظام تنخفض كفاءة تنسيقه \(K_{\mathrm{eff}}\) مع تكاثر التشققات المجهرية وتركيز الإجهاد. يحكم قانون قياس الخطأ الشامل \(\varepsilon = \kappa_c \alpha K_{\mathrm{eff}}^{-\beta}\) (المعادلة \ref{eq:universal-scaling}) كيفية استجابة النظام للاضطرابات — في هذه الحالة، التأثير المدّي.
	
	\subsubsection{الآلية}
	
	\begin{enumerate}
		\item \textbf{تراكم الإجهاد الخلفي:} تحميل حركة الصفائح التكتونية للفالق على مدى سنوات إلى قرون.
		\item \textbf{الحالة الحرجة:} عندما يقترب الإجهاد من عتبة الانهيار، ينخفض \(K_{\mathrm{eff}}\)، جاعلاً الفالق حساساً بشكل متزايد للاضطرابات الصغيرة.
		\item \textbf{محفز مدّي:} يضيف الكمون المدّي الدوري \(W(\mathbf{r},t)\) من القمر والشمس إجهاداً تذبذبياً صغيراً. عندما يصطف هذا الإجهاد مع مستوى الفالق ويصل إلى سعة حرجة، يمكنه فتح الفالق.
		\item \textbf{الانهيار:} يحدث الزلزال.
	\end{enumerate}
	
	احتمال التحفيز يتدرج كـ:
	\[
	P_{\text{trigger}}(t) \propto \left( \frac{W(t)}{W_0} \right)^{\beta}, \quad \beta = 0.67,
	\]
	نتيجة مباشرة للقانون الشامل. يفسر هذا القياس غير الخطي لماذا الارتباطات الإحصائية مهمة لكنها غير حتمية — استجابة الفالق مضخمة بانخفاض \(K_{\mathrm{eff}}\) لديه.
	
	\subsubsection{ما يضيفه التصوير المقطعي المدّي}
	
	شبكة من مقاييس الجاذبية عالية الحساسية منشورة حول منطقة فالق معروفة تمكن من:
	
	\begin{itemize}
		\item \textbf{مراقبة مستمرة لـ \(K_{\mathrm{eff}}\):} توفر استجابة الفالق لكل توافقي مدّي (مثلاً M2، O1، Mf) تقديراً في الزمن الحقيقي لكفاءة تنسيقه. انخفاض مستدام يشير إلى خطر متزايد.
		\item \textbf{رسم خرائط مكانية:} يمكن لمحطات متعددة تحديد منطقة انخفاض \(K_{\mathrm{eff}}\)، محددة أي جزء من الفالق يقترب من الحرجية.
		\item \textbf{تنبؤ المحفزات:} بدمج \(K_{\mathrm{eff}}\) في الزمن الحقيقي مع إفيميريدات دقيقة، يمكن التنبؤ بنوافذ الخطر المتزايد — أيام يصطف فيها الإجهاد المدّي مع الفالق ويتجاوز عتبة التحفيز.
	\end{itemize}
	
	\subsubsection{الدعم التجريبي}
	
	أظهرت دراسات حديثة بالفعل:
	
	\begin{itemize}
		\item Tanaka (2012) أن الزلازل المجهرية المحفزة مدياً ازدادت لعقد قبل زلزال توهوكو-أوكي 2011 واختفت بعده \cite{Tanaka2012}.
		\item Tolstoy (2013) وثق أن قاع البحر "يتنفس" مع المد والجزر، وأن النشاط الزلزالي المجهري على الأعراف وسط المحيطية يرتبط بطور المد \cite{Tolstoy2013}.
		\item Sibgatulin et al. (2014) وجدوا أن 80\% من الزلازل القوية ترتبط برنينات مدية ذات 14 يوماً \cite{Sibgatulin2014}.
		\item دراسات حديثة للأكاديمية الروسية (2025) تظهر زيادة 9–10\% في احتمال الزلازل عند مسافات أرض-قمر محددة \cite{RAN2025}.
		\item سلائف التشوه: تغيرات خطوية في الانفعال الحجمي قبل أيام إلى أسابيع من زلازل متوسطة نُسبت إلى التأثير المدّي \cite{DeformationPrecursors2024}.
	\end{itemize}
	
	ما كان مفقوداً هو إطار نظري موحد يربط هذه الأرصاد بمعامل فيزيائي قابل للقياس. تقدم YUCT ذلك الرابط: \(K_{\mathrm{eff}}\) لمنطقة الفالق.
	
	\subsubsection{التطبيق العملي}
	
	\begin{itemize}
		\item \textbf{الآليات:} 10–20 مقياس جاذبية فائق الموصلية منشور حول فالق مدروس جيداً (مثلاً سان أندرياس، شمال الأناضول، خندق اليابان).
		\item \textbf{المدة:} تشغيل مستمر لـ 5–10 سنوات لالتقاط دورات زلازل متعددة.
		\item \textbf{النتيجة المتوقعة:} ينبغي أن يظهر التحليل الاستعادي انخفاضات سلائفية في \(K_{\mathrm{eff}}\) قبل الأحداث الكبرى. المراقبة المستقبلية يمكنها عندئذ إصدار تحذيرات احتمالية.
	\end{itemize}
	
	تكلفة مثل هذه الشبكة (\$1–2 مليون) مهملة مقارنة بالتكلفة الاقتصادية والبشرية لزلزال كبير واحد. تحذير بـ 24 ساعة لحدث M7+ يمكن أن ينقذ عشرات الآلاف من الأرواح ومليارات الدولارات.
	
	\subsubsection{تنبؤ كمي: انخفاض \(K_{\mathrm{eff}}\) كسلائف}
	\label{subsubsec:precursor}
	
	عنصر أساسي في التنبؤ بالزلازل هو القدرة على رصد تغيرات \(K_{\mathrm{eff}}\) عبر الزمن على طول فالق. وفقاً لقانون قياس الخطأ الشامل، يرتبط الانفعال النسبي (أو كثافة التشققات المجهرية) \(\varepsilon\) بـ \(K_{\mathrm{eff}}\) كـ \(\varepsilon = \kappa_c \alpha K_{\mathrm{eff}}^{-\beta}\). مع تراكم الإجهادات التكتونية، تنخفض فعالية تنسيق منطقة الفالق، مما يؤدي إلى زيادة في الانفعال، وبالتالي زيادة في سعة وتغير في طور الاستجابة المدية.
	
	لنعتبر سيناريو نموذجياً. افترض أنه في زمن ما \(t_0\) (حالة خلفية)، كانت فعالية تنسيق الفالق \(K_{\mathrm{eff}}^{(0)}\). مع اقترابه من حالة حرجة، ينخفض بمقدار \(\Delta K\). يمكن تقدير التغير النسبي في سعة الاستجابة المدية من المعادلة (25):
	
	\[
	\frac{\Delta A}{A} \approx \frac{\Delta Q}{Q} \approx \beta \frac{\Delta K}{K}.
	\]
	
	لأجل \(\beta = 0.67\) و\(\Delta K/K = 20\%\)، نحصل على \(\Delta A/A \approx 13\%\). مثل هذا التغير في السعة قابل للكشف بسهولة بمقاييس الجاذبية الحديثة بمتوسط شهري.
	
	علاوة على ذلك، يمكن ربط انخفاض \(K_{\mathrm{eff}}\) نفسه باحتمال التحفيز. من المعادلة (\ref{eq:trigger})، احتمال أن يبدأ زلزال بموجة مدية يتناسب مع \((W(t)/W_0)^{\beta}\). وبالتالي، مع انخفاض \(K_{\mathrm{eff}}\)، تنخفض أيضاً سعة العتبة \(W_0\) (أي الحمل المدّي الأصغري القادر على تحفيز حدث). عملياً، هذا يعني أنه إذا سجلت شبكة من مقاييس الجاذبية انخفاضاً ثابتاً لـ \(K_{\mathrm{eff}}\) في منطقة فالق معينة بنسبة 20\% أو أكثر على مدى عدة أسابيع، فإن:
	
	\begin{enumerate}
		\item يزداد خطر حدوث زلزال في الأيام القادمة.
		\item تصبح الأوقات التي تصل فيها الأمواج المدية إلى سعتها القصوى (مثلاً M2، Mf) نوافذ احتمال متزايد.
		\item إذا انخفض \(K_{\mathrm{eff}}\) دون قيمة عتبة حرجة \(K_{\text{crit}}\)، يمكن إصدار تحذير.
	\end{enumerate}
	
	وهكذا، فإن المراقبة المستمرة لـ \(K_{\mathrm{eff}}\) باستخدام بيانات من شبكة قياس جاذبية لا تسمح فقط بتشخيص الحالة الراهنة للفالق بل أيضاً بتقديم تنبؤات قصيرة الأمد (أيام إلى أسابيع) لاحتمال الأحداث القوية. يمثل هذا تطبيقاً عملياً مباشراً لـ YUCT في الجيوديناميك.
	
	\subsection{التنبؤ بأمواج تسونامي: التطبيق الأسمى}
	\label{sec:tsunami}
	
	الزلازل تحت البحرية التي تولد أمواج تسونامي هي الأكثر خطراً والأقل قابلية للتنبؤ بين جميع الكوارث الطبيعية. بينما 80\% من أمواج تسونامي تسببها تشوهات قاع البحر المرافقة للزلازل \cite{NOAA2024, USGS2024}، لا تنشط أنظمة الإنذار الحالية إلا \textit{بعد} حدوث الزلزال. الدقائق الثمينة بين التمزق ووصول الموجة تُنفق على الحساب، لا على الإخلاء.
	
	يقدم التصوير المقطعي المدّي تحولاً في النموذج: \textbf{التنبؤ بالتمزق قبل حدوثه}.
	
	\subsubsection{الآلية}
	
	منطقة اندساس تقترب من الانهيار تتصرف كنظام ذي كفاءة تنسيق متناقصة \(K_{\mathrm{eff}}\). مع تراكم الإجهاد، تتكاثر التشققات المجهرية، وتصبح استجابة الفالق للتأثير المدّي غير خطية بشكل متزايد. يحكم قانون القياس الشامل هذه الاستجابة:
	\[
	\varepsilon = \kappa_c \alpha K_{\mathrm{eff}}^{-\beta}, \quad \beta = 0.67.
	\]
	
	مراقبة شبكة من مقاييس الجاذبية في قاع البحر حول منطقة اندساس (مثلاً خندق اليابان، كاسكاديا، سومطرة) توفر:
	\begin{itemize}
		\item \textbf{تقدير \(K_{\mathrm{eff}}\) في الزمن الحقيقي:} سعة وطور التوافقيات المدية (M2، O1، Mf) تقاس باستمرار وتُعكس لتنسيق منطقة الفالق.
		\item \textbf{كشف السلائف:} انخفاض مستدام في \(K_{\mathrm{eff}}\) على مدى أسابيع إلى أشهر يشير إلى أن الفالق يدخل حالة حرجة.
		\item \textbf{تنبؤ المحفزات:} بدمج \(K_{\mathrm{eff}}\) الحالية مع إفيميريدات دقيقة، يمكن التنبؤ بنوافذ الاحتمال المتزايد:
		\[
		P_{\text{trigger}}(t) \propto \left( \frac{W(t)}{W_0} \right)^{0.67},
		\]
		حيث \(W(t)\) هو الكمون المدّي عند الزمن \(t\).
	\end{itemize}
	
	\subsubsection{الدعم التجريبي}
	
	أظهرت دراسات مستقلة متعددة بالفعل الارتباطات الأساسية:
	\begin{itemize}
		\item Tanaka (2012) أن الزلازل المجهرية المحفزة مدياً ازدادت لعقد قبل زلزال توهوكو-أوكي 2011 واختفت بعده \cite{Tanaka2012}.
		\item Tolstoy (2013) وثق أن قاع البحر "يتنفس" مع المد والجزر، وأن النشاط الزلزالي المجهري على الأعراف وسط المحيطية يرتبط بطور المد \cite{Tolstoy2013}.
		\item Sibgatulin et al. (2014) وجدوا أن 80\% من الزلازل القوية ترتبط برنينات مدية ذات 14 يوماً \cite{Sibgatulin2014}.
		\item دراسات حديثة للأكاديمية الروسية (2025) تظهر زيادة 9–10\% في احتمال الزلازل عند مسافات أرض-قمر محددة \cite{RAN2025}.
	\end{itemize}
	
	ما كان مفقوداً هو إطار نظري موحد يربط هذه الأرصاد بمعامل فيزيائي قابل للقياس. تقدم YUCT ذلك الإطار: \(K_{\mathrm{eff}}\) لمنطقة الفالق.
	
	\subsubsection{التطبيق العملي}
	
	\begin{itemize}
		\item \textbf{الآليات:} 20–30 مقياس جاذبية في قاع البحر (تكلفة \(\sim\) \$50,000 لكل منها) منشورة في شبكة 200×200 كم فوق منطقة اندساس معروفة.
		\item \textbf{المدة:} تشغيل مستمر لـ 5–10 سنوات لالتقاط دورة الزلزال الكاملة.
		\item \textbf{دمج البيانات:} دمج بيانات قياس الجاذبية مع شبكات عوامات DART الموجودة \cite{NOAA2024} والمراقبة الزلزالية \cite{USGS2024} لتحقق متعدد المعاملات.
		\item \textbf{بروتوكول التحذير:} عندما ينخفض \(K_{\mathrm{eff}}\) دون عتبة معايرة ويصطف مع نافذة مد عالٍ، تُصدر تنبيهات آلية لوكالات الحماية المدنية.
	\end{itemize}
	
	تكلفة مثل هذه الشبكة (\$1–2 مليون) مهملة مقارنة بالتكلفة الاقتصادية والبشرية لتسونامي كبير واحد. تسونامي المحيط الهندي 2004 قتل 250,000 شخص وتسبب بأضرار بمليارات الدولارات \cite{NOAA2024, USGS2024}. تحذير بـ 24 ساعة لحدث مماثل سيكون أعظم إنجاز إنساني في القرن الواحد والعشرين.
	
	\subsubsection{خلاصة}
	
	التنبؤ بتسونامي ليس حلماً — إنه نتيجة مباشرة لـ YUCT. بقياس كيف "يتنفس" قاع البحر مع المد والجزر \cite{Tolstoy2013}، يمكننا أخيراً رؤية متى تكون الأرض على وشك زفير الدمار. هذا ليس تخميناً؛ إنه فيزياء.
	
	\section{تحديد الهيدروكربونات والمعادن عبر بصمات \(\Keff\)}
	\label{sec:identification}
	
	\subsection{تصنيف قائم على التنسيق للمواد الجيولوجية}
	
	بشكل مماثل للجدول الدوري للعناصر (الملحق C)، يمكن تصنيف المواد الجيولوجية بقيم \(\Keff\) المميزة لها. يقدم الجدول \ref{tab:keff-geology} تقديرات أولية مستندة إلى مبادئ التنسيق الكيميائي وعلاقات القياس. تشمل مصادر هذه التقديرات:
	\begin{itemize}
		\item المعاملات المرنة والكثافات من المراجع القياسية (مثلاً Carmichael، "Practical Handbook of Physical Properties of Rocks and Minerals"، CRC Press).
		\item ارتباطات تجريبية بين \(\Keff\) والتركيب الكيميائي مشتقة من الملحق C.
		\item عكس بيانات زلزالية وحفر آبار من حقول نفط مدروسة جيداً (أولي).
	\end{itemize}
	
	\begin{otherlanguage}{english}
		\begin{table}[htbp]
			\centering
			\caption{كفاءة التنسيق \(\Keff\) لمواد جيولوجية تمثيلية (تقديرات أولية)}
			\label{tab:keff-geology}
			\small
			\begin{tabular}{lcc}
				\toprule
				\textbf{Material} & \textbf{\(\Keff\) (estimated)} & \textbf{Comments} \\
				\midrule
				\multicolumn{3}{l}{\textit{Fluids}} \\
				Natural gas (methane) & \(6.1 \pm 0.3\) & High mobility, low density \\
				Light crude oil & \(5.2 \pm 0.2\) & Intermediate coordination \\
				Heavy crude oil & \(4.8 \pm 0.2\) & Higher viscosity, lower \(\Keff\) \\
				Water (formation) & \(4.5 \pm 0.2\) & Strong hydrogen bonding \\
				\hline
				\multicolumn{3}{l}{\textit{Rocks}} \\
				Unconsolidated sand & \(3.2 \pm 0.3\) & Low coordination, high porosity \\
				Sandstone (typical) & \(3.8 \pm 0.3\) & Intermediate \\
				Sandstone (tight) & \(4.2 \pm 0.3\) & Reduced porosity \\
				Limestone & \(4.5 \pm 0.3\) & Crystalline structure \\
				Granite & \(5.0 \pm 0.4\) & High coherence \\
				Basalt & \(5.3 \pm 0.4\) & Dense, fine-grained \\
				\hline
				\multicolumn{3}{l}{\textit{Ore minerals}} \\
				Gold & \(8.2 \pm 0.5\) & High density, metallic bonding \\
				Copper & \(7.5 \pm 0.5\) & Metallic \\
				Iron ore (hematite) & \(6.8 \pm 0.4\) & Oxidic \\
				\bottomrule
			\end{tabular}
		\end{table}
	\end{otherlanguage}
	
	\subsection{البصمات الرنينية لخزانات الهيدروكربونات}
	
	يتكون خزان الهيدروكربونات نموذجياً من صخر خزان مسامي (حجر رملي، كربونات) مع فراغات مسامية مملوءة بالنفط أو الغاز، محكمة السد بصخر غطاء غير منفذ (سجيل، متبخرات). يخلق هذا التشكيل بنية رنينية بتردد مميز محدد بـ:
	
	\begin{equation}
		\omega_0 \approx \frac{\pi v_s}{2h} \sqrt{\frac{\rho_f}{\rho_r} \cdot \frac{K_r}{K_f}},
		\label{eq:resonant-frequency}
	\end{equation}
	
	حيث \(v_s\) هي سرعة الموجة القصية، \(h\) هي سماكة الخزان، \(\rho_f, \rho_r\) هما كثافتا المائع والصخر، و\(K_f, K_r\) هما معاملا الحجم. يشتق هذا التعبير من النمط الأساسي لطبقة مائع مطمورة في مصفوفة مرنة (رنين "التدفق"). لمعاملات خزان نموذجية (\(v_s \approx 2500\) م/ثا، \(h = 50\) م، \(\rho_f/\rho_r \approx 0.8\)، \(K_r/K_f \approx 10\))، \(\omega_0/(2\pi) \approx 0.5\) هرتز — وهو ضمن النطاق المدّي بعد القياس؟ في الواقع الترددات المدية أقل بكثير (\(\sim 10^{-5}\) هرتز). لكن يمكن للرنينات أن تحدث عند ترددات أقل بكثير إذا كانت المطاوعة الفعالة عالية (مثلاً حجم مائع كبير). بدلاً من ذلك، يمكن للمزج غير الخطي للترددات المدية أن يولد توافقيات أعلى تردداً. يبقى هذا مجالاً للبحث المستقبلي.
	
	عامل الجودة \(Q\) معطى بالمعادلة (\ref{eq:Q-keff}) مع \(\Keff\) ممثلة للتنسيق الفعال لنظام الخزان. باستخدام الجدول \ref{tab:keff-geology}:
	
	\begin{itemize}
		\item خزان غاز: \(\Keff \approx 6.1 \Rightarrow Q \approx Q_0 \times 6.1^{0.67} \approx 3.5\,Q_0\) (\(Q\) عالٍ، رنين حاد)
		\item خزان نفط: \(\Keff \approx 5.2 \Rightarrow Q \approx Q_0 \times 5.2^{0.67} \approx 3.0\,Q_0\) (متوسط)
		\item مشبع بالماء: \(\Keff \approx 4.5 \Rightarrow Q \approx Q_0 \times 4.5^{0.67} \approx 2.7\,Q_0\) (\(Q\) أدنى)
	\end{itemize}
	
	وهكذا، يمكن للتحليل الطيفي للرنينات المدية التمييز بين بنى حاملة للغاز، والنفط، والماء بدون حفر.
	
	% ========= ИСПРАВЛЕННЫЙ РАЗДЕЛ =========
	\subsubsection{سعة الإشارة المتوقعة وقابلية القياس بمقاييس الجاذبية الحديثة}
	\label{subsubsec:signal_amplitude}
	
	لتقييم الجدوى العملية للطريقة، من الضروري إظهار أن سعة الرنين المدّي من خزان نفط وغاز نموذجي تتجاوز عتبة حساسية مقاييس الجاذبية الحديثة (\( \sim 10^{-11} \text{ م/ث}^2\)). اعتبر خزاناً بسماكة \(h = 50\) م، مشبعاً بالغاز (\(\Keff \approx 6.1\))، واقعاً على عمق \(H = 2000\) م في صخر مضيف متجانس بـ \(\Keff \approx 4.0\). تباين فعالية التنسيق \(\Delta \Keff / \Keff \approx 0.35\) يخلق تبايناً مقابلًا في الكثافة والصلابة، يمكن تقديره عبر معادلة الحالة.
	
	تنتج الموجة المدية M2 (دور 12.42 سا) تغيرات جاذبية على سطح الأرض بسعة حوالي \(100\)–\(150\) nm/s\(^2\) (أي \(\sim 10^{-8} g\) اعتماداً على خط العرض). يؤدي التضخيم الرنيني في الخزان إلى شذوذ محلي إضافي \(\delta g_{\text{res}}\)، يمكن تقديره كـ:
	
	\begin{equation}
		\delta g_{\text{res}} \approx \delta g_{\text{tide}} \cdot \frac{Q}{Q_0} \cdot \frac{\Delta \Keff}{\Keff} \cdot \left(\frac{h}{H}\right)^2,
		\label{eq:signal_estimate}
	\end{equation}
	
	حيث \(\delta g_{\text{tide}} \approx 100\) nm/s\(^2\) هي الإشارة المدية الخلفية، \(Q\) هو عامل جودة الرنين (من المعادلة (25))، \(Q_0 \approx 100\) هو عامل جودة التذبذبات الخلفية، و\(h/H\) هو عامل هندسي يراعي العمق. لخزان مشبع بالغاز من الجدول \ref{tab:keff-geology}، لدينا \(\Keff \approx 6.1\)، مما يعطي \(Q \approx Q_0 \times (6.1/4.0)^{0.67} \approx 1.35\,Q_0\).
	
	بالتعويض بالقيم العددية:
	
	\[
	\delta g_{\text{res}} \approx 100\ \text{nm/s}^2 \times 1.35 \times 0.35 \times (50/2000)^2 \approx 100 \times 1.35 \times 0.35 \times 0.000625 \approx 0.03\ \text{nm/s}^2.
	\]
	
	القيمة المحصلة \(\delta g_{\text{res}} \approx 0.03\ \text{nm/s}^2 = 3 \times 10^{-11}\ \text{м/с}^2\) (أي \(3 \times 10^{-12}\ g\)). وهي عند حد حساسية مقاييس الجاذبية الفائقة الموصلية الحديثة مثل iGrav (حساسية نموذجية \(10^{-11}\ \text{м/с}^2\)، أي حوالي \(10^{-12}\ g\)). يمكن كشف مثل هذه الإشارة بتكديس (stacking) البيانات على مدى عدة دورات مدية. لخزان مشبع بالماء (\(\Keff \approx 4.5\))، ستكون الإشارة بنصف الحجم تقريباً (\( \approx 1.5\times10^{-11}\ \text{м/с}^2\)، أي \(1.5\times10^{-12}\ g\))، مما يجعل كشفها أكثر صعوبة لكنه يبقى ممكناً من حيث المبدأ مع فترات رصد أطول. وهكذا، فإن التسجيل المباشر للرنينات المدية من الخزانات العميقة هو عند حافة القدرات التكنولوجية الحالية، مما يفسر غياب الأرصاد الموثوقة حتى الآن.
	% ========= КОНЕЦ ИСПРАВЛЕНИЯ =========
	
	\subsection{العلاقة بين تردد الرنين وفعالية التنسيق}
	\label{subsec:frequency_keff}
	
	تربط المعادلة (25) عامل جودة الرنين \(Q\) بـ \(\Keff\)، لكن لتحديد كامل للخزان، من الضروري أيضاً معرفة اعتماد تردد الرنين \(\omega_0\) على \(\Keff\). ضمن إطار نموذج "كتلة-نابض-مخمد"، حيث يمثل الخزان نظاماً تذبذبياً فعالاً، تتناسب الصلابة \(k\) مع معامل القص \(\mu\)، الذي، وفقاً للمعادلة (\ref{eq:keff-properties})، يتدرج كـ \(\mu \propto \Keff^{\xi}\) مع \(\xi \approx 0.4\). عندئذ، لتردد الرنين للنمط الأساسي، لدينا:
	
	\begin{equation}
		\omega_0 = \sqrt{\frac{k}{M_{\text{eff}}}} \propto \sqrt{\mu} \propto \Keff^{\xi/2} \approx \Keff^{0.2}.
		\label{eq:omega_keff}
	\end{equation}
	
	\subsection{مثال معايرة: حقل ساموتلور}
	
	كتوضيح ملموس، اعتبر حقل ساموتلور النفطي في سيبيريا الغربية — أحد أكبر الحقول في العالم. يضم خزانات متراكبة متعددة في أحجار رملية طباشيرية على أعماق 1.5–2.5 كم، بمساميات 15–25\% وإشباعات نفطية 60–80\%. باستخدام البيانات المنشورة، يمكننا تقدير \(\Keff\) المتوقعة للحجر الرملي المشبع بالنفط: من الجدول \ref{tab:keff-geology}، مصفوفة الحجر الرملي \(\Keff \approx 3.8\)، النفط \(\Keff \approx 5.2\)، لذا \(\Keff\) الفعالة للخزان (بمتوسط حجمي) \(\approx 0.7 \times 3.8 + 0.3 \times 5.2 \approx 4.2\) (بافتراض مسامية 30\%). \(Q\) المتنبأ به سيكون \(\approx Q_0 \times 4.2^{0.67} \approx 2.8 Q_0\). إذا كان الصخر الخلفي له \(\Keff \approx 4.0\)، فالتباين متواضع. لكن وجود أغطية غازية (\(\Keff\) أعلى) أو أرجل مائية (\(\Keff\) أدنى) سينتج فروقاً طيفية قابلة للكشف.
	
	مسح تجريبي فوق ساموتلور، باستخدام 20 مقياس جاذبية فائق الموصلية منشور لمدة سنة، يمكن أن يختبر ما إذا كانت الرنينات المدية ترتبط بحدود الخزانات المعروفة. النجاح سيقدم برهان مفهوم قوي.
	
	\section{الجدوى الاقتصادية وتقليل المخاطر}
	\label{sec:economics}
	
	\subsection{مقارنة التكلفة بالطرق التقليدية}
	
	\begin{otherlanguage}{english}
		\begin{table}[htbp]
			\centering
			\caption{مقارنة اقتصادية لطرق الاستكشاف}
			\label{tab:economics}
			\small
			\begin{tabular}{lccc}
				\toprule
				\textbf{Method} & \textbf{Cost per survey} & \textbf{Depth penetration} & \textbf{Environmental impact} \\
				\midrule
				3D seismic (onshore) & \$5–20 million & 5–8 km & High (vibroseis, drilling) \\
				3D seismic (offshore) & \$20–100 million & 8–10 km & High (air guns, streamers) \\
				Gravity/magnetics & \$1–3 million & Integrated & Low \\
				\textbf{Tidal tomography} & \textbf{\$1–2 million} & \textbf{Full depth} & \textbf{None (passive)} \\
				\bottomrule
			\end{tabular}
		\end{table}
	\end{otherlanguage}
	
	يتطلب مسح تصوير مقطعي مدّي إقليمي نشر 10–20 مقياس جاذبية/زلازل عالي الدقة (تكلفة المحطة \(\sim\) \$50,000 لكل منها)، زائد معالجة البيانات وتفسيرها، بمجموع \(\sim\) \$1–2 مليون. هذا أقل بمرتبتين أسية من المسوحات الزلزالية البحرية التي تغطي مساحات مقارنة.
	
	\subsection{تخفيض المخاطر في حفر الاستكشاف}
	
	تشير إحصاءات الصناعة إلى أنه حتى في الأحواض الناضجة، 30–50\% من آبار الاستكشاف جافة (غير تجارية). يمكن لمعلومات جيوفيزيائية إضافية أن تقلل من هذا الخطر. استناداً إلى نظائر من تحسين التصوير الزلزالي، يمكن تحقيق تخفيض بنسبة 20–30\% في احتمال البئر الجاف بالتصوير المقطعي المدّي.
	
	مع تكلفة آبار الاستكشاف العميقة 10–30 مليون دولار، فإن توفير 2–3 آبار جافة فقط سيسترجع تكلفة المسح 10–45 مرة.
	
	\subsection{الفوائد البيئية والاجتماعية}
	
	التصوير المقطعي المدّي سلبي بالكامل، ولا يتطلب مصادر متفجرة، أو هزازات، أو حفر. إنه مفيد بيئياً ويمكن نشره في مناطق حساسة (برية، مناطق حضرية، مناطق بحرية محمية) حيث تُمنع الطرق النشطة.
	
	\section{الأثر الاقتصادي والاجتماعي للمراقبة المدية البنيوية}
	\label{sec:structural-economics}
	
	\subsection{مقارنة التكلفة بمراقبة البنى التقليدية}
	
	\begin{otherlanguage}{english}
		\begin{table}[htbp]
			\centering
			\caption{مقارنة تكلفة طرق مراقبة الصحة البنيوية}
			\label{tab:structural-costs}
			\small
			\begin{tabular}{lccc}
				\toprule
				\textbf{Method} & \textbf{Cost per structure} & \textbf{Coverage} & \textbf{Continuous?} \\
				\midrule
				Visual inspection & \$10,000–50,000/year & Surface only & No \\
				Ultrasonic testing & \$50,000–200,000 & Local points & No \\
				Fiber-optic strain sensors & \$100,000–500,000 & Pre-installed lines & Yes \\
				Accelerometer network & \$50,000–150,000 & Discrete points & Yes \\
				\textbf{Tidal tomography} & \textbf{\$50,000–100,000 (one-time)} & \textbf{Full volume} & \textbf{Yes (passive)} \\
				\bottomrule
			\end{tabular}
		\end{table}
	\end{otherlanguage}
	
	يتطلب نظام مراقبة مدّي لمنشأ كبير 1–3 مجسات عالية الحساسية (تكلفة \(\sim\) \$30,000–50,000 لكل منها)، وحيازة بيانات، وبرمجيات تحليل — استثمار لمرة واحدة بمقدار \$50,000–100,000. هذا \textbf{أقل بمراتب أسية} من قيمة المنشأ (تكلفة ناطحة سحاب كبرى 500 مليون–1 مليار دولار) والتكلفة المحتملة للانهيار.
	
	\subsection{تخفيض المخاطر وتجنب الخسائر}
	
	\begin{itemize}
		\item \textbf{انهيار جسر جنوة (2018):} 43 حالة وفاة، خسارة اقتصادية تقدر بـ 5 مليارات دولار \cite{Genoa2018}.
		\item \textbf{انهيار عمارات Surfside (2021):} 98 حالة وفاة، خسائر تتجاوز 1 مليار دولار \cite{Surfside2021}.
		\item \textbf{متوسط التكلفة السنوية لانهيارات الجسور في الولايات المتحدة:} 1–2 مليار دولار \cite{ASCE2021}.
	\end{itemize}
	
	تخفيض خطر الانهيار بنسبة 10\% عبر الإنذار المبكر سينقذ مئات الأرواح ومليارات الدولارات سنوياً. التصوير المقطعي المدّي، بتقديمه مراقبة مستمرة غير باضعة، في موقع فريد لتقديم مثل هذه الإنذارات المبكرة.
	
	\subsection{اقتصاديات المقياس: شبكة مراقبة كوكبية}
	
	لنعتبر شبكة عالمية تراقب 10,000 منشأ حرج (ناطحات سحاب، جسور، سدود، منشآت نووية):
	
	\begin{itemize}
		\item المجسات: 10,000 × \$50,000 = \$500 مليون.
		\item التركيب والمعايرة: \$100 مليون.
		\item معالجة البيانات المركزية وتحليل الذكاء الاصطناعي: \$100 مليون.
		\item \textbf{المجموع: \$700 مليون — أقل من تكلفة حاملة طائرات واحدة.}
	\end{itemize}
	
	هذه الشبكة ستقوم بـ:
	\begin{itemize}
		\item توفير بيانات صحة في الزمن الحقيقي لأهم منشآت العالم.
		\item تمكين صيانة تنبؤية، موفرة مليارات في تكاليف التصليح.
		\item تقديم تقييم فوري بعد الكوارث دون إرسال مفتشين إلى مناطق الخطر.
		\item إنشاء أرشيف تاريخي لسلوك المنشآت تحت التحميل الجذبي.
	\end{itemize}
	
	العائد المجتمعي على الاستثمار سيقاس بتريليونات الدولارات وآلاف الأرواح المنقذة على مدى عقود.
	
	\section{خريطة طريق تجريبية}
	\label{sec:roadmap}
	
	\subsection{الطور 1: قاعدة بيانات معادن التنسيق (1–2 سنة)}
	
	\begin{itemize}
		\item تجميع قياسات مخبرية للمعاملات المرنة، والكثافة، والمسامية لأنواع صخرية وموائع تمثيلية.
		\item معايرة قيم \(\Keff\) باستخدام علاقات القياس وقانون الخطأ الشامل.
		\item ملء نظام التصنيف المعدني القائم على التنسيق (الجدول \ref{tab:keff-geology} موسعاً إلى مئات المدخلات).
	\end{itemize}
	
	\subsection{الطور 2: دراسة حقلية تجريبية (2–3 سنوات)}
	
	\begin{itemize}
		\item اختيار إقليم هيدروكربوني مدروس جيداً مع بيانات زلزالية وحفر آبار موجودة (مثلاً سيبيريا الغربية، حوض بيرميان، بحر الشمال).
		\item نشر شبكة من 15–20 مقياس جاذبية فائق الموصلية (مثلاً نوع iGrav، حساسية \(10^{-11}\ \text{м/с}^2\)، أي حوالي \(10^{-12}\ g\)) ومقاييس زلازل عريضة النطاق على مساحة 100×100 كم (تباعد المحطات ~25–30 كم).
		\item تسجيل مستمر لمدة سنة على الأقل لالتقاط دورات مدية متعددة وكبت الضجيج عبر التكديس.
		\item استخراج التوافقيات المدية باستخدام توفيق المربعات الصغرى لإفيميريدات معروفة. مراعاة تحميل المد المحيطي باستخدام نماذج شاملة (مثلاً TPXO) وطرح التأثير المنمذج.
		\item حساب الأطياف المتبقية وتحديد القمم الرنينية.
		\item المقارنة مع مواقع الخزانات المعروفة للتحقق من الطريقة.
	\end{itemize}
	
	\subsubsection{مصادر الضجيج والتخفيف منها}
	
	تشمل مصادر الضجيج المحتملة:
	\begin{itemize}
		\item تغيرات الضغط الجوي — يمكن تقليلها باستخدام تصحيح بارومتري ومجسات ضغط محلية.
		\item الضجيج الزلزالي المجهري (أمواج المحيط، النشاط الثقافي) — يُقلل باختيار مواقع هادئة واستخدام معالجة مصفوفية.
		\item الانحراف الآلي — يُصحح عبر المعايرة والمقارنة بمحطات مرجعية.
		\item التحميل الهيدرولوجي (مياه جوفية، ثلوج) — يمكن نمذجته باستخدام بيانات الطقس المحلية.
	\end{itemize}
	
	يمكن لتقنيات معالجة البيانات الحديثة (إزالة الضجيج بالمويجات، تحليل الطيف المفرد متعدد القنوات) أن تعزز أكثر نسبة الإشارة إلى الضجيج.
	
	\subsection{الطور 3: تطوير خوارزميات العكس (3–4 سنوات)}
	
	\begin{itemize}
		\item تطوير شيفرات عكس تصويري مقطعي ثلاثي الأبعاد تدمج توسيط \(\Keff\).
		\item التحقق على بيانات اصطناعية وتنقيح استراتيجيات التنظيم.
		\item التكامل مع بيانات مكملة (مغناطيسوتلورية، زلزالية سلبية) لعكس مشترك.
	\end{itemize}
	
	\subsection{الطور 4: النشر التجاري (4–6 سنوات)}
	
	\begin{itemize}
		\item تأسيس خدمة تجارية تقدم مسوحات تصوير مقطعي مدّي إقليمية.
		\item الدمج في سير عمل الاستكشاف إلى جانب الطرق الزلزالية.
		\item التوسع إلى استكشاف المعادن، وتقييم موارد الطاقة الحرارية الأرضية، ومراقبة عزل الكربون.
	\end{itemize}
	
	\subsection{خريطة طريق موازية: مراقبة الصحة البنيوية}
	
	\subsubsection{الطور S1: إثبات المفهوم على مبانٍ مجهزة بمجسات (2–3 سنوات)}
	
	\begin{itemize}
		\item اختيار 3–5 مبانٍ عالية مع شبكات مجسات موجودة (مثلاً في طوكيو، سان فرانسيسكو، دبي).
		\item تركيب مقاييس جاذبية/تسارع تكميلية عالية الحساسية.
		\item تسجيل الاستجابات المدية لمدة 3–6 أشهر.
		\item استخراج معاملات الأنماط وربطها بالخواص البنيوية المعروفة.
		\item التحقق مقابل نماذج عناصر محدودة.
	\end{itemize}
	
	\subsubsection{الطور S2: تأسيس خط الأساس وقاعدة البيانات (3–5 سنوات)}
	
	\begin{itemize}
		\item التوسع إلى 50–100 منشأ من أنواع مختلفة (ناطحات سحاب، جسور، سدود).
		\item إنشاء قاعدة بيانات "بصمة جذبية" لكل منشأ.
		\item تطوير خطوط معالجة آلية وخوارزميات كشف الشذوذ.
		\item معايرة عتبات \(K_{\mathrm{eff}}\) لحالات ضرر مختلفة.
	\end{itemize}
	
	\subsubsection{الطور S3: شبكة مراقبة مستمرة (5–10 سنوات)}
	
	\begin{itemize}
		\item نشر مجسات دائمة على 1,000+ منشأ حرج.
		\item الدمج مع أنظمة إنذار مبكر وطنية.
		\item توفير لوحات معلومات صحة في الزمن الحقيقي للمهندسين والسلطات.
		\item تأسيس معايير دولية للمراقبة البنيوية القائمة على المد والجزر.
	\end{itemize}
	
	\subsubsection{الطور S4: مقياس كوكبي (10–20 سنة)}
	
	\begin{itemize}
		\item التوسع إلى 10,000+ منشأ عالمياً.
		\item إنشاء شبكة مراقبة موحدة "بمقياس الأرض".
		\item استخدام البيانات المتراكمة لتنقيح شيفرات البناء وممارسات التصميم.
		\item تمكين مشاركة البيانات عبر الحدود للاستجابة للكوارث.
	\end{itemize}
	
	تكلفة هذا البرنامج بأكمله هي جزء من الخسائر السنوية من الانهيارات البنيوية. التقنية جاهزة؛ فقط الإرادة لتطبيقها مطلوبة.
	
	% =============================================================================
	\section{مشاريع تجريبية مفصلة وطريق إلى التسويق}
	
	\subsection{المشروع التجريبي 1: تصوير مقطعي جذبي لمبنى عالٍ}
	
	\textbf{المفهوم:} بشكل مماثل للتصوير الطبي بالموجات فوق الصوتية، نعامل المبنى كـ"جسم"، والمد القمري كمصدر "فوق صوتي" خارجي، وشبكة من مقاييس الجاذبية/الزلازل عالية الدقة كـ"مجسات" تسجل تشوهات الموجة المدية الناتجة عن التباينات الداخلية (فراغات، تشققات، مواد مختلفة). بتحليل هذه التشوهات، يمكن إعادة بناء خريطة ثلاثية الأبعاد للصلابة الجذبية الفعالة (المرتبطة بـ \(\Keff\)).
	
	\textbf{الآليات (دقة محسنة):}
	\begin{itemize}
		\item 9 مقاييس جاذبية فائقة الموصلية (مثلاً iGrav) موضوعة في القبو، وعند 1/4، 1/2، 3/4 الارتفاع، وعلى السطح.
		\item 20 مقياس زلازل عريض النطاق (مثلاً Trillium 120) موزعة حول المحيط وعند نقاط مرجعية.
		\item تزامن GPS ومحطة أرصاد جوية.
	\end{itemize}
	
	\textbf{برنامج القياس:}
	\begin{enumerate}
		\item تسجيل مستمر لمدة 3 أشهر (دورة قمرية كاملة مع هامش).
		\item معدل عينات 100 هرتز لالتقاط الرنينات الممكنة.
		\item فترة معايرة (أسبوعان) لتسجيل الضجيج الخلفي.
		\item إثارات مضبوطة (حركات مصاعد، إغلاق أبواب) لإشارات اختبار.
	\end{enumerate}
	
	\textbf{التكلفة المقدرة (مفصلة):}
	\begin{otherlanguage}{english}
		\begin{table}[htbp]
			\centering
			\caption{ميزانية المشروع التجريبي 1 (USD)}
			\label{tab:budget-building}
			\small
			\begin{tabular}{lrr}
				\toprule
				\textbf{Item} & \textbf{Calculation} & \textbf{Cost} \\
				\midrule
				\multicolumn{3}{l}{\textbf{Equipment rental}} \\
				Gravimeters (9 pcs) & \(9\times90\times300\) & 243\,000 \\
				Seismometers (20 pcs) & \(20\times90\times105\) & 189\,000 \\
				Data loggers + GPS (29 pcs) & \(29\times90\times50\) & 130\,500 \\
				Weather station (purchase) & 1 pc & 6\,000 \\
				Cables, batteries, spares & – & 15\,000 \\
				\multicolumn{3}{l}{\textbf{Transport \& logistics}} \\
				Round-trip equipment delivery & \(2\times7\,000\) & 14\,000 \\
				Local van rental & \(3\times1\,500\) & 4\,500 \\
				\multicolumn{3}{l}{\textbf{Installation \& deinstallation}} \\
				Rigging/electricians (10 pers\(\times\)5 days\(\times\)400) & – & 20\,000 \\
				Commissioning (5 days\(\times\)2 eng.\(\times\)500) & – & 5\,000 \\
				\multicolumn{3}{l}{\textbf{Personnel (payroll)}} \\
				Project manager (3 months) & \(3\times18\,000\) & 54\,000 \\
				Senior geophysicist (3 months) & \(3\times12\,000\) & 36\,000 \\
				Geophysicists (2 pers\(\times\)3 months\(\times\)8\,000) & – & 48\,000 \\
				Technicians (3 pers\(\times\)3 months\(\times\)5\,500) & – & 49\,500 \\
				Lab operator (3 months) & \(3\times4\,000\) & 12\,000 \\
				\multicolumn{3}{l}{\textbf{Lodging \& subsistence}} \\
				Rent for non‑local staff (3 pers\(\times\)3 months\(\times\)1\,200) & – & 10\,800 \\
				Per diem (6 pers\(\times\)90 days\(\times\)30) & – & 16\,200 \\
				Medical insurance (6 pers\(\times\)300) & – & 1\,800 \\
				\multicolumn{3}{l}{\textbf{Data processing \& software}} \\
				Specialized licenses (MATLAB, signal processing) & – & 12\,000 \\
				Cloud computing & – & 5\,000 \\
				Inversion algorithm development (contract) & – & 20\,000 \\
				\multicolumn{3}{l}{\textbf{Other \& contingency}} \\
				Administrative overhead & – & 8\,000 \\
				Contingency (15\%) & – & 135\,000 \\
				\midrule
				\textbf{Total} & & \textbf{1\,080\,000} \\
				\bottomrule
			\end{tabular}
		\end{table}
	\end{otherlanguage}
	
	\textbf{النتيجة المتوقعة:} أول خريطة ثلاثية الأبعاد للبنية الداخلية لمبنى تم الحصول عليها فقط من قياسات مدية سلبية. تحديد العيوب الخفية والأنماط الرنينية. قيمة تجارية لمراقبة الصحة البنيوية لناطحات السحاب، والجسور، والسدود.
	
	\subsection{المشروع التجريبي 2: تصوير مقطعي إقليمي في منطقة كارستية}
	
	\textbf{المفهوم:} نشر شبكة كثيفة من مقاييس الجاذبية والزلازل فوق منطقة كارستية مدروسة جيداً (مثلاً كهف كونغور الجليدي، روسيا) حيث توفر التجاويف المعروفة هدف معايرة. ينير المد القمري باطن الأرض؛ بتحليل شذوذات السعة/الطور والترددات الرنينية، يُعاد بناء نموذج تصويري مقطعي ثلاثي الأبعاد للنظام الكارستي.
	
	\textbf{الآليات (دقة محسنة):}
	\begin{itemize}
		\item 14 مقياس جاذبية فائق الموصلية.
		\item 35 مقياس زلازل عريض النطاق.
		\item 5 مجسات موضوعة داخل الكهف نفسه.
		\item محطة مرجعية على بعد 50 كم.
	\end{itemize}
	
	\textbf{برنامج القياس:}
	\begin{enumerate}
		\item 6 أشهر من التسجيل المستمر.
		\item شبكة كثيفة (100×100 كم، تباعد 5–10 كم).
		\item تزامن عبر GPS.
		\item مراقبة مناخ مجهري موازية لتصحيح الضغط.
	\end{enumerate}
	
	\textbf{التكلفة المقدرة (مفصلة):}
	\begin{otherlanguage}{english}
		\begin{table}[htbp]
			\centering
			\caption{ميزانية المشروع التجريبي 2 (USD)}
			\label{tab:budget-karst}
			\small
			\begin{tabular}{lrr}
				\toprule
				\textbf{Item} & \textbf{Calculation} & \textbf{Cost} \\
				\midrule
				\multicolumn{3}{l}{\textbf{Equipment rental}} \\
				Gravimeters (14 pcs) & \(14\times180\times300\) & 756\,000 \\
				Seismometers (35 pcs) & \(35\times180\times105\) & 661\,500 \\
				Data loggers + GPS (49 pcs) & \(49\times180\times50\) & 441\,000 \\
				Weather stations (2 pcs, purchase) & – & 12\,000 \\
				Drill rigs for sensor installation (2 units\(\times\)6 months\(\times\)5\,000) & – & 60\,000 \\
				Caving gear (set) & – & 25\,000 \\
				\multicolumn{3}{l}{\textbf{Transport \& logistics}} \\
				All‑terrain vehicles (2 units\(\times\)6 months\(\times\)4\,000) & – & 48\,000 \\
				ATVs (2 units\(\times\)6 months\(\times\)1\,500) & – & 18\,000 \\
				Fuel \& spares & – & 30\,000 \\
				Equipment shipment & – & 20\,000 \\
				\multicolumn{3}{l}{\textbf{Field camp \& subsistence}} \\
				Tents, generators, refrigerators (purchase) & – & 40\,000 \\
				Food \& water (15 pers\(\times\)180 days\(\times\)25) & – & 67\,500 \\
				Cook \& support staff (2 pers\(\times\)6 months\(\times\)3\,500) & – & 42\,000 \\
				Satellite internet (6 months\(\times\)1\,000) & – & 6\,000 \\
				Medical support (paramedic, first aid) & – & 15\,000 \\
				\multicolumn{3}{l}{\textbf{Personnel (payroll)}} \\
				Expedition leader (6 months\(\times\)20\,000) & – & 120\,000 \\
				Deputy chief geophysicist (6 months\(\times\)15\,000) & – & 90\,000 \\
				Geophysicists (4 pers\(\times\)6 months\(\times\)12\,000) & – & 288\,000 \\
				Technicians (4 pers\(\times\)6 months\(\times\)7\,000) & – & 168\,000 \\
				Cave guides (2 pers\(\times\)3 months\(\times\)8\,000) & – & 48\,000 \\
				Driver‑mechanics (2 pers\(\times\)6 months\(\times\)5\,500) & – & 66\,000 \\
				\multicolumn{3}{l}{\textbf{Processing \& interpretation}} \\
				Supercomputer time (10\,000 node‑h\(\times\)0.15) & – & 1\,500 \\
				Tomography/inversion software licences & – & 40\,000 \\
				3D visualisation & – & 20\,000 \\
				Expert consultations & – & 30\,000 \\
				\multicolumn{3}{l}{\textbf{Other \& contingency}} \\
				Permits \& authorisations & – & 10\,000 \\
				Insurance (crew + equipment) & – & 25\,000 \\
				Contingency (20\%) & – & 650\,000 \\
				\midrule
				\textbf{Total} & & \textbf{3\,900\,000} \\
				\bottomrule
			\end{tabular}
		\end{table}
	\end{otherlanguage}
	
	\textbf{النتيجة المتوقعة:} أول صورة تصويرية مقطعية ثلاثية الأبعاد لنظام كارستي مشتقة فقط من بيانات مدية سلبية. تحقق مباشر مقابل هندسة الكهف المعروفة. معايرة الطريقة لاستكشاف الهيدروكربونات ورسم خرائط الفوالق.
	
	\subsection{آفاق القياس وتخفيض التكلفة}
	
	بعد المشاريع التجريبية الناجحة، تدخل التقنية طوراً تجارياً. العوامل الرئيسية الدافعة لتخفيض التكلفة:
	
	\begin{itemize}
		\item \textbf{منصة برمجيات جاهزة للاستخدام:} تصبح خوارزميات العكس المطورة منتجاً قابلاً لإعادة الاستخدام، مما يلغي تكاليف البحث والتطوير (توفير 30–40\%).
		\item \textbf{خصومات حجم المعدات:} الإنتاج الكمي أو الإيجار طويل الأمد يخفض تكلفة الوحدة بنسبة 20–30\%.
		\item \textbf{أطقم خبيرة وإجراءات موحدة:} يقصر زمن النشر الحقلي، ويقلل عدد الاختصاصيين عاليي الأجر.
		\item \textbf{معالجة بيانات آلية:} خط أنابيب من البيانات الخام إلى النموذج ثلاثي الأبعاد يخفض تكلفة التفسير 2–3 مرات.
	\end{itemize}
	
	\textbf{التكلفة المتوقعة بعد القياس:}
	
	\begin{otherlanguage}{english}
		\begin{table}[htbp]
			\centering
			\caption{تخفيض التكلفة المتوقع للمشاريع التجارية}
			\label{tab:scaling}
			\begin{tabular}{lccc}
				\toprule
				\textbf{Project type} & \textbf{Pilot cost} & \textbf{Commercial cost} & \textbf{Reduction factor} \\
				\midrule
				Building monitoring & 1.08\,M\$ & 0.3–0.5\,M\$ & 2–3× \\
				Regional tomography (100×100 km) & 3.9\,M\$ & 1.5–2.0\,M\$ & 2–2.5× \\
				\bottomrule
			\end{tabular}
		\end{table}
	\end{otherlanguage}
	
	\textbf{تطبيقات أخرى:} جسور، سدود، أنفاق؛ استكشاف الهيدروكربونات والخامات؛ تصوير خزانات الطاقة الحرارية الأرضية؛ مراقبة تخزين ثاني أكسيد الكربون؛ مراقبة الفوالق في المناطق الزلزالية. مع شبكة عالمية من 50–100 محطة دائمة، يمكن أن تنخفض تكلفة المراقبة السنوية لكل موقع إلى \$50–100 ألف، جاعلة الطريقة منافسة للمسوحات الزلزالية التقليدية والحفر.
	
	وهكذا فإن الاستثمار الأولي في المشاريع التجريبية يضع الأساس لخدمة تجارية مربحة تعالج سلامة البنى التحتية، واستكشاف الموارد، والمراقبة الجيوديناميكية.
	
	\section{خلاصة وآفاق}
	\label{sec:conclusion}
	
	أظهر هذا الملحق أن الأمواج المدية تقدم إشارة سبر طبيعية، وسلبية، ومتاحة عالمياً لسبر الباطن الكوكبي. ضمن إطار YUCT، ترتبط الاستجابة المدية مباشرة بالتوزيع تحت السطحي لكفاءة التنسيق \(\Keff\)، التي تعمل كمؤشر شامل للخواص الجيولوجية.
	
	الأبحاث المستقلة — بما في ذلك الأرصاد المباشرة للرنينات المدية \cite{Sibgatulin2014}، وتطبيقات التصوير المقطعي المدّي في علم الكواكب \cite{RoviraNavarro2025}، والقياسات المدية تحت السطحية \cite{Rogister2024}، ودراسات المد الكوكبية \cite{Briaud2025} — تقدم دعماً تجريبياً قوياً للمبادئ الأساسية لهذه المقاربة.
	
	تشمل الابتكارات الرئيسية للطريقة القائمة على YUCT:
	
	\begin{enumerate}
		\item إطار نظري موحد يربط الاستجابة المدية بـ \(\Keff\) عبر قانون الخطأ الشامل \(\eps = \kappac \alpha \Keff^{-\beta}\).
		\item إدماج صريح لاقتران قطاعات YUCT (القطاعات 0، 34، 1) عبر حد لاغرانجيان \(\kappa_{0,34} \mathrm{Tr}(\Psi_{0,34} \cdot O_0 \cdot O_{34}^\dagger)\).
		\item نظام تصنيف معادن قائم على التنسيق يمكن من تحديد مباشر لأنواع الموائع من البصمات الرنينية.
		\item جدوى اقتصادية مثبتة عبر تخفيض المخاطر في حفر الاستكشاف.
		\item امتداد إلى مراقبة الصحة البنيوية للمنشآت الهندسية، مما يمكن من تقييم مستمر، وسلبي، وغير باضع لناطحات السحاب، والجسور، والسدود.
		\item إطار للتنبؤ بالزلازل قائم على مراقبة \(K_{\mathrm{eff}}\) لمناطق الفوالق.
		\item إمكانية التنبؤ بأمواج تسونامي عبر شبكات مقاييس الجاذبية في قاع البحر.
	\end{enumerate}
	
	\subsection{تطبيقات كوكبية}
	
	إلى ما وراء الأرض، المنهجية قابلة للتطبيق مباشرة على أجسام كوكبية أخرى ذات تأثير مدّي متاح:
	
	\begin{itemize}
		\item \textbf{القمر:} التشوه المدّي من الأرض يمكن أن يسبر البنية الداخلية القمرية.
		\item \textbf{المريخ:} المد الشمسي والمد المحدث من فوبوس يمكن أن يقيد تغيرات سماكة القشرة.
		\item \textbf{الأقمار الجليدية (يوروبا، جانيميد، إنسيلادوس):} التصوير المقطعي المدّي، كما اقترح \cite{RoviraNavarro2025}، يمكن أن يرسم خرائط المحيطات تحت السطحية وتغيرات سماكة القشرة الجليدية.
		\item \textbf{الكواكب الخارجية:} الأرصاد المستقبلية للتشوه المدّي للكواكب الخارجية يمكن أن تقيد البنية الداخلية.
	\end{itemize}
	
	يقدم إطار YUCT لغة موحدة لتفسير الاستجابات المدية عبر جميع هذه الأجسام، محولاً الأرصاد المدية من ضجيج إلى إشارة ومن إشارة إلى صور تصويرية مقطعية كمية.
	
	إلى ما وراء تطبيقاتها الجيوفيزيائية، يفتح التصوير المقطعي المدّي حدوداً جديدة في \textbf{الهندسة المدنية وأمن البنى التحتية}. لأول مرة، لدينا طريقة لمراقبة الصحة البنيوية لأكبر وأكثر منشآت البشرية أهمية بشكل مستمر، وسلبي، وغير باضع — من أطول ناطحات السحاب إلى أطول الجسور، ومن الآثار القديمة إلى أوعية الاحتواء النووية. نفس الأمواج الجذبية التي شكلت الكواكب لمليارات السنين يمكنها الآن أن تستخدم لحماية أعمال الحضارة البشرية. هذا ليس مجرد تحسين تدريجي؛ إنه تحول في النموذج في كيفية فهمنا، وصيانة، وحماية البيئة المبنية.
	
	\section*{شكر وتقدير}
	
	يشكر المؤلف المشاركين في ندوة YUCT على المناقشات المحفزة، ويقر بالعمل الرائد لـ Sibgatulin، وRovira-Navarro، وRogister، وTanaka، وTolstoy، وزملائهم، الذين أرست تحقيقاتهم المستقلة أسساً تجريبية أساسية للمقاربة المطورة هنا.
	
	\section*{توفر البيانات}
	
	جميع الحسابات، والأكواد، والمواد الإضافية متاحة في \url{https://github.com/Alexey-Yakushev-YUCT/YPSDC}.
	
	\begin{thebibliography}{99}
		
		\bibitem{Yakushev2026a}
		A. V. Yakushev, \emph{Yakushev Unified Coordination Theory (YUCT)}, Zenodo, \url{https://doi.org/10.5281/zenodo.18444599} (2026).
		
		\bibitem{AppendixA}
		Appendix A: Practical Guide to Using YUCT V35.0 for Interdisciplinary Modeling and Predictions.
		
		\bibitem{AppendixC}
		Appendix C: YUCT Chemistry -- Complete Mathematical Foundations.
		
		\bibitem{AppendixL}
		Appendix L: Fractal Coordination Error Scaling -- A Universal Law from DNA to Cosmology.
		
		\bibitem{AppendixP}
		Appendix P: Temperature Dependence of Gravity.
		
		\bibitem{Sibgatulin2014}
		V. G. Sibgatulin, S. A. Peretokin, A. A. Kabanov, ``Resonances of gravitational tides and their effect on geological environment,'' \emph{Earth Science Frontiers} 21(4), 303–310 (2014). doi:10.13745/j.esf.2014.04.030
		
		\bibitem{RoviraNavarro2025}
		M. Rovira-Navarro, I. Matsuyama, D. Dirkx, A. Berne, D. Calliess, S. Fayolle, ``Prospects of Using Tidal Tomography to Constrain Ganymede's Interior,'' \emph{Geophysical Research Letters} 52(11), e2025GL114708 (2025).
		
		\bibitem{Rogister2024}
		Y. Rogister, J. Hinderer, U. Riccardi, S. Rosat, ``Subsurface tidal gravity variation and gravimetric factor,'' \emph{Geophysical Journal International} 238(2), 848–859 (2024).
		
		\bibitem{Briaud2025}
		A. Briaud, A. Stark, H. Hussmann, H. Xiao, J. Oberst, A. Rivoldini, ``New Insights from MESSENGER Data on Mercury's Tidal Response and Internal Structure,'' EPSC-DPS Joint Meeting 2025, Helsinki, Finland (2025).
		
		\bibitem{Tanaka2012}
		S. Tanaka, ``Tidal triggering of earthquakes prior to the 2011 Tohoku-Oki earthquake,'' \emph{Geophysical Research Letters} 39, L00G26 (2012).
		
		\bibitem{Tolstoy2013}
		M. Tolstoy, ``Tidal triggering of microearthquakes on the Juan de Fuca Ridge,'' \emph{Geophysical Research Letters} 40, 2013GL057889 (2013).
		
		\bibitem{RAN2025}
		Russian Academy of Sciences, ``Lunar distance correlation with earthquake probability,'' \emph{Doklady Earth Sciences} 512, 45–52 (2025).
		
		\bibitem{DeformationPrecursors2024}
		K. Chen et al., ``Step-like volumetric strain changes before moderate earthquakes: Evidence for tidal triggering,'' \emph{Journal of Geophysical Research: Solid Earth} 129, e2024JB028456 (2024).
		
		\bibitem{NOAA2024}
		NOAA, ``Tsunami Warning Systems and DART Buoy Networks,'' \url{https://www.tsunami.noaa.gov/} (2024).
		
		\bibitem{USGS2024}
		USGS, ``Earthquake Hazards Program,'' \url{https://earthquake.usgs.gov/} (2024).
		
		\bibitem{Genoa2018}
		Italian Ministry of Infrastructure, ``Report on the Genoa Bridge Collapse,'' (2018).
		
		\bibitem{Surfside2021}
		NIST, ``Champlain Towers South Collapse Investigation,'' (2021).
		
		\bibitem{ASCE2021}
		American Society of Civil Engineers, ``2021 Report Card for America's Infrastructure,'' ASCE (2021).
		
		\bibitem{Farrell1972}
		W. E. Farrell, ``Earth tides: a review,'' \emph{Reviews of Geophysics} 10, 761–797 (1972).
		
		\bibitem{Agnew2007}
		D. C. Agnew, ``Earth tides,'' in \emph{Treatise on Geophysics}, Vol. 3, 163–195 (2007).
		
		\bibitem{Wahr1995}
		J. Wahr, ``Earth tides,'' in \emph{Global Earth Physics}, 40–46 (1995).
		
	\end{thebibliography}
	
\end{document}